\documentclass{article}
\usepackage{amsmath, amssymb, amsthm}
\usepackage{geometry}
\usepackage[UTF8]{ctex}

\newtheorem{problem}{问题}
\newtheorem{solution}{解答}
\title{数值分析第十一周作业}
\date{\today}
\author{李佩优PB23000270}
\geometry{a4paper,scale=0.8}

\begin{document}
\maketitle

\section*{8.5}

\subsection*{1}
讨论下列多步法的稳定性和相容性:

根据定理1，多步法收敛的充要条件是稳定且相容。将各方法写成标准形式：
\[
a_k x_n + a_{k-1}x_{n-1} + \cdots + a_0 x_{n-k} = h\left(b_k f_n + b_{k-1}f_{n-1} + \cdots + b_0 f_{n-k}\right)
\]
定义特征多项式 \(p(\lambda)=\sum_{j=0}^k a_j\lambda^j\) 和 \(q(\lambda)=\sum_{j=0}^k b_j\lambda^j\)。

稳定性：\(p(\lambda)\) 的所有根满足 \(|\lambda|\le 1\)，且模为1的根都是单根。 

相容性：\(p(1)=0\) 且 \(p'(1)=q(1)\)。

a. \(x_n - x_{n-2} = 2h f_{n-1}\)

取 \(k=2\)，有 \(a_2=1,\ a_1=0,\ a_0=-1\)；\(b_2=0,\ b_1=2,\ b_0=0\)。  
\(p(\lambda)=\lambda^2-1\)，根为 \(\lambda=1,-1\)，均为模为1的单根 → 稳定。  
\(p(1)=0\)，\(p'(\lambda)=2\lambda,\ p'(1)=2\)；\(q(\lambda)=2\lambda,\ q(1)=2\) → 相容。  

$\Rightarrow$收敛。

b. \(x_n - x_{n-2} = h\left[\frac{7}{3}f_{n-1} - \frac{2}{3}f_n + \frac{1}{3}f_{n-3}\right]\)

最大下标偏移为3，取 \(k=3\)。写为标准形式：
\[
x_n + 0\cdot x_{n-1} - x_{n-2} + 0\cdot x_{n-3} = h\left(-\frac{2}{3}f_n + \frac{7}{3}f_{n-1} + 0\cdot f_{n-2} + \frac{1}{3}f_{n-3}\right)
\]
即 \(a_3=1,\ a_2=0,\ a_1=-1,\ a_0=0\)；\(b_3=-\frac{2}{3},\ b_2=\frac{7}{3},\ b_1=0,\ b_0=\frac{1}{3}\)。  
\(p(\lambda)=\lambda^3-\lambda=\lambda(\lambda-1)(\lambda+1)\)，根为 \(0,1,-1\)，模为1的根都是单根 → 稳定。  
\(p(1)=0\)，\(p'(\lambda)=3\lambda^2-1,\ p'(1)=2\)；\(q(\lambda)=-\frac{2}{3}\lambda^3+\frac{7}{3}\lambda^2+\frac{1}{3}\)，\(q(1)=-\frac{2}{3}+\frac{7}{3}+\frac{1}{3}=2\) → 相容。  

$\Rightarrow$收敛。

c. \(x_n - x_{n-1} = h\left[\frac{3}{8}f_n + \frac{19}{24}f_{n-1} - \frac{5}{24}f_{n-2} + \frac{1}{24}f_{n-3}\right]\)

取 \(k=3\)，标准形式：
\[
x_n - x_{n-1} + 0\cdot x_{n-2} + 0\cdot x_{n-3} = h\left(\frac{3}{8}f_n + \frac{19}{24}f_{n-1} - \frac{5}{24}f_{n-2} + \frac{1}{24}f_{n-3}\right)
\]
即 \(a_3=1,\ a_2=-1,\ a_1=0,\ a_0=0\)；\(b_3=\frac{3}{8},\ b_2=\frac{19}{24},\ b_1=-\frac{5}{24},\ b_0=\frac{1}{24}\)。  
\(p(\lambda)=\lambda^3-\lambda^2=\lambda^2(\lambda-1)\)，根为 \(\lambda=0\)（二重），\(\lambda=1\)（单根）。模为1的根是单根 → 稳定。  
\(p(1)=0\)，\(p'(\lambda)=3\lambda^2-2\lambda,\ p'(1)=1\)；\(q(1)=\frac{3}{8}+\frac{19}{24}-\frac{5}{24}+\frac{1}{24}=1\) → 相容。  

$\Rightarrow$收敛。

\subsection*{6}
判断下列多步法中哪一个是收敛的：

a. \(x_n - x_{n-2} = h(f_n - 3f_{n-1} + 4f_{n-2})\)

取 \(k=2\)：\(a_2=1, a_1=0, a_0=-1\)；\(b_2=1, b_1=-3, b_0=4\)。  
\(p(\lambda)=\lambda^2-1\)，根为 \(\pm1\)（单根）→ 稳定。  
\(p(1)=0\)，\(p'(1)=2\)；\(q(\lambda)=\lambda^2-3\lambda+4\)，\(q(1)=1-3+4=2\) → 相容。  

$\Rightarrow$收敛。

b. \(x_n - 2x_{n-1} + x_{n-2} = h(f_n - f_{n-1})\)

\(a_2=1, a_1=-2, a_0=1\)；\(p(\lambda)=\lambda^2-2\lambda+1=(\lambda-1)^2\)。根 \(\lambda=1\) 是重根，不稳定。  

$\Rightarrow$不收敛。

c. \(x_n - x_{n-1} - x_{n-2} = h(f_n - f_{n-1})\)

\(a_2=1, a_1=-1, a_0=-1\)；\(p(\lambda)=\lambda^2-\lambda-1\)，根为 \(\frac{1\pm\sqrt{5}}{2}\)，其中 \(\frac{1+\sqrt{5}}{2}>1\)，不稳定。  

$\Rightarrow$不收敛。

d. \(x_n - 3x_{n-1} + 2x_{n-2} = h(f_n + f_{n-1})\)

\(a_2=1, a_1=-3, a_0=2\)；\(p(\lambda)=\lambda^2-3\lambda+2=(\lambda-1)(\lambda-2)\)，根 \(\lambda=2>1\)，不稳定。  

$\Rightarrow$不收敛。

e. \(x_n - x_{n-2} = h(f_n - 3f_{n-1} + 2f_{n-2})\)

\(a_2=1, a_1=0, a_0=-1\)；\(p(\lambda)=\lambda^2-1\)，稳定。但相容性：  
\(p'(1)=2\)；\(q(\lambda)=\lambda^2-3\lambda+2\)，\(q(1)=1-3+2=0 \neq 2\)，不相容。  

$\Rightarrow$不收敛。

\section*{8.12}
\subsection*{4}


首先将多步法写成标准形式：
\[
x_n + \alpha x_{n-1} - (1+\alpha)x_{n-2} = h\!\left[-\frac{a}{2}f_n + \frac{4+3\alpha}{2}f_{n-1}\right],
\]
其中 \(f_k = f(t_k, x_k)\)。对应生成多项式为
\[
\rho(r) = r^2 + \alpha r - (1+\alpha),\qquad
\sigma(r) = -\frac{a}{2}r^2 + \frac{4+3\alpha}{2}r.
\]

1. 相容性：

方法相容要求 \(\rho(1)=0\)（已满足）且 \(\rho'(1)=\sigma(1)\)。
\[
\rho'(1) = 2+\alpha,\qquad \sigma(1) = \frac{4+3\alpha - a}{2}.
\]
令两者相等得 \(4+2\alpha = 4+3\alpha - a\)，即 \(a = \alpha\)。原方法变为
\[
x_n + \alpha x_{n-1} - (1+\alpha)x_{n-2} = \frac{h}{2}\bigl[-\alpha f_n + (4+3\alpha)f_{n-1}\bigr].
\]

2. 零稳定性：

\(\rho(r) = r^2 + \alpha r - (1+\alpha) = (r-1)(r+1+\alpha)\)。零稳定要求所有根满足 \(|r|\le 1\)，且单位圆上的根均为单根。根为 \(r_1=1\)，\(r_2=-(1+\alpha)\)。条件：
\[
|1+\alpha| \le 1 \quad\text{且}\quad \alpha \neq -2\; (\text{否则 }r_2=1\text{ 为重根}).
\]
解得 \(-2 < \alpha \le 0\)。

3. 收敛性：

零稳定与相容性同时满足即收敛，故 \(\alpha\in(-2,0]\)。

4. 阶（至少二阶）：

计算局部截断误差，将真解展开代入可得误差主项为
\[
\frac{4+5\alpha}{12}h^3 x^{(3)}(t_n) + O(h^4).
\]
方法对任意 \(\alpha\) 至少为二阶；当 \(\alpha \neq -4/5\) 时恰好二阶。所求的 A-稳定区间内 \(\alpha \neq -0.8\)，故均满足二阶要求。

5. A-稳定性：

将方法用于试验方程 \(x' = \lambda x\)，令 \(z = h\lambda\)，得特征方程
\[
\left(1 + \frac{\alpha}{2}z\right)r^2 + \left[\alpha - \frac{4+3\alpha}{2}z\right]r - (1+\alpha) = 0.
\]
要求对任意 \(\mathrm{Re}(z)<0\) 均有 \(|r|<1\)。
通过双线性变换 \(r = \frac{1+w}{1-w}\) 或直接计算边界 \(\mu(\theta) = \rho(e^{i\theta})/\sigma(e^{i\theta})\) 的实部：
\[
\mathrm{Re}\,\mu(\theta) = \frac{4\alpha(\alpha+1)(1-\cos\theta)^2}{|\sigma(e^{i\theta})|^2}.
\]
A-稳定要求边界不进入左半平面，即 \(\mathrm{Re}\,\mu(\theta) \ge 0\)，故 \(\alpha(\alpha+1) \ge 0\)，得 \(\alpha \le -1\) 或 \(\alpha \ge 0\)。另外当 \(z\to -\infty\) 时，两根趋于 \(0\) 和 \(\frac{4+3\alpha}{\alpha}\)，必须满足 \(\left|\frac{4+3\alpha}{\alpha}\right| \le 1\)，等价于 \(\alpha \in [-2,-1]\)。结合零稳定，得到 A-稳定的范围是 \(\alpha \in (-2,-1]\)。（当 \(\alpha \ge 0\) 时不满足无穷远条件；\(\alpha=0\) 为中点公式，非 A-稳定。）

\(\alpha\) 的取值范围为 \(-2 < \alpha \le -1\)。


\subsection*{6}

隐式梯形法则为
\[
x_n - x_{n-1} = \frac{h}{2}\bigl[f_n + f_{n-1}\bigr].
\]
应用于试验方程 \(x' = \lambda x\)，令 \(z = h\lambda\)，得
\[
x_n - x_{n-1} = \frac{z}{2}(x_n + x_{n-1})
\;\Longrightarrow\;
\left(1 - \frac{z}{2}\right)x_n = \left(1 + \frac{z}{2}\right)x_{n-1}.
\]
增长因子 \(R(z) = \dfrac{1 + z/2}{1 - z/2}\)。绝对稳定性要求 \(|R(z)| < 1\)。
设 \(z = x + \mathrm{i}y\)，则
\[
|1 + z/2|^2 < |1 - z/2|^2
\;\Longrightarrow\;
(1 + x/2)^2 + (y/2)^2 < (1 - x/2)^2 + (y/2)^2.
\]
化简得 \(x < 0\)。

因此隐式梯形法则的绝对稳定性区域为左半开平面 \(\{z \in \mathbb{C} : \mathrm{Re}(z) < 0\}\)。

\end{document} 